\documentclass[10pt]{beamer}
\usetheme{metropolis}
\usepackage{graphicx}
\usepackage{booktabs}
\usepackage{tabularx}
\usepackage{amsmath}
\usepackage{amssymb}
\usepackage{bm}
\usepackage{mathtools}
\usepackage[style=numeric,sorting=none,maxbibnames=3,minbibnames=1]{biblatex}
\addbibresource{lit.bib}

% Custom color overrides (white background, black text/titles)
\setbeamercolor{background canvas}{bg=white}
\setbeamercolor{normal text}{fg=black}
\setbeamercolor{frametitle}{bg=white, fg=black}
\setbeamercolor{section title}{fg=black}
\setbeamercolor{palette primary}{bg=white, fg=black}

% Custom font overrides (make slide titles bigger)
\setbeamerfont{frametitle}{size=\Large}

% Custom BME red/burgundy styling for elements
\definecolor{BMEred}{RGB}{128, 0, 0}
\setbeamercolor{alerted text}{fg=BMEred}
\setbeamercolor{example text}{fg=BMEred}

\title{Chaos in Large Language Models}
\subtitle{Thesis Defense}
\author{Jaca Gregorio}
\date{2026}
\institute{Budapest University of Technology and Economics (BME)}

\begin{document}

% Slide 1: BME Title Page (matches thesis title page layout)
\begin{frame}[plain]
  \centering
  \rule{\linewidth}{0.5mm} \\[0.3cm]
  {\Large \textbf{Chaos in Large Language Models}} \\[0.1cm]
  \rule{\linewidth}{0.5mm} \\[0.4cm]
  
  \begin{columns}[t]
    \begin{column}{0.48\textwidth}
      \begin{flushleft}
        \small \textbf{Author:} \\
        \textbf{Jaca Gregorio} \\
        {\scriptsize Physicist Engineering BSc.}
      \end{flushleft}
    \end{column}
    \begin{column}{0.48\textwidth}
      \begin{flushleft}
        \small \textbf{Supervisors:} \\
        \textbf{Dr. Török János} \\
        {\scriptsize Associate Professor} \\[0.05cm]
        \textbf{Benedek Kristóf} \\
        {\scriptsize PhD Candidate} \\[0.05cm]
        {\tiny BME Theoretical Physics Department}
      \end{flushleft}
    \end{column}
  \end{columns}
  
  \vfill
  \includegraphics[scale=0.4]{plots/bme.png} \\[0.1cm]
  {\small \textbf{BME} \\ 2026}
\end{frame}

% \section{Introduction and Motivation}

% Slide 2: Introduction and Motivation
\begin{frame}{Introduction and Motivation}
  LLMs are highly non-linear systems. They are Deep Neural Networks trained on large amounts of data.
  
  \vspace{0.5cm}
  
  There is growing interest in interpretability and understanding the model's dynamics.
  
  \vspace{0.5cm}
  
  Reproducibility and determinism is important in order to do Reinforcement Learning on-policy. Quantifying sensitivity to perturbations and initial conditions is necessary.
  
  \vspace{0.5cm}
  
  They are an interesting case study for chaos theory.
\end{frame}



% Slide 5: Nonlinear Dynamics and Chaos
\begin{frame}{Nonlinear Dynamics and Chaos}
  \textbf{Deterministic Map Evolution:}
  \begin{equation*}
    \mathbf{x}_{t+1} = \mathbf{f}(\mathbf{x}_t)
  \end{equation*}
  
  \textbf{Sensitivity to Initial Conditions:}
  \begin{equation*}
    \|\bm{\delta}(t)\| \approx \|\bm{\delta}(0)\| e^{\lambda t}
  \end{equation*}
  \begin{itemize}
    \item $\lambda > 0$: Maximal Lyapunov Exponent (strong indicator of chaos).
    \item Exponential divergence makes long-term prediction impossible.
  \end{itemize}
  
  \textbf{Stretching and Folding:}
  \begin{itemize}
    \item Dissipative systems converge to a bounded \textbf{strange attractor}.
    \item Stretching (stretching trajectories) + Folding (confinement to bounds) creates a fractal attractor structure of non-integer correlation dimension $D_2$.
  \end{itemize}
\end{frame}

% Slide 6: Recurrence Plots
\begin{frame}{Recurrence Plots}
  \begin{columns}
    \begin{column}{0.35\textwidth}
      % \textbf{Mathematical Formulation:}
      \begin{equation*}
        R_{i,j}(\epsilon) = \Theta(\epsilon - d(\mathbf{x}_i, \mathbf{x}_j))
      \end{equation*}
      \begin{itemize}
        \item $\epsilon$: recurrence threshold.
        \item $\Theta$: Heaviside step function.
      \end{itemize}
      \textbf{Qualitative Types:}
      \begin{itemize}
        \item Periodic: Long parallel diagonals.
        \item Random: Uniform noise.
        \item Chaotic: Interrupted diagonals.
      \end{itemize}
    \end{column}
    \begin{column}{0.63\textwidth}
      \centering
      \includegraphics[width=0.48\textwidth]{figures_pdf/RP/compare_systems/sine_rp.pdf}
      \includegraphics[width=0.48\textwidth]{figures_pdf/RP/compare_systems/lorenz_36_rp.pdf} \\
      \includegraphics[width=0.48\textwidth]{figures_pdf/RP/compare_systems/white_noise_rp.pdf}
      \includegraphics[width=0.48\textwidth]{figures_pdf/RP/compare_systems/brownian_rp_0.pdf}
      \\[0.1cm]
      {\tiny Fig 1: Recurrence plots of (a) periodic, (b) chaotic Lorenz, (c) white noise, (d) Brownian motion.}
    \end{column}
  \end{columns}
\end{frame}






% Slide 8: The Chaos-LLM Analogy
\begin{frame}{The Chaos-LLM Analogy}
  \centering
  \renewcommand{\arraystretch}{1.2}
  \footnotesize
  \begin{tabular}{l l}
    \toprule
    \textbf{Chaotic System Concept} & \textbf{LLM Equivalent} \\ 
    \midrule
    Phase Space & Embedding / Hidden State Space \\
    State Vector & Hidden State Vector $\mathbf{x}$ \\
    Trajectory & Token sequence (Horizontal) / Layer evolution (Vertical) \\
    Time Step & Token index $t$ (Horizontal) / Layer index $\ell$ (Vertical) \\
    Evolution Map & Autoregressive generation / Layer-wise updates \\
    Nonlinear Coupling & Self-Attention Mechanism \\
    Stretching & Attention / MLP Blocks \\
    Folding & Normalization Layers \\
    Initial Condition & Input prompt / Initial layer hidden state \\
    Measurement Error & Floating-Point Precision / Quantization \\ 
    \bottomrule
  \end{tabular}
\end{frame}



% Slide 10: Dimensionality & Power Iteration
\begin{frame}{Dimensionality \& Jacobian Power Iteration}
  \textbf{Correlation Dimension ($D_2$) Estimation:}
  \begin{itemize}
    \item Derived from correlation integral $C(\epsilon) \sim \epsilon^{D_2}$:
      \begin{equation*}
        C(\epsilon) = \lim_{N_{\text{traj}} \to \infty} \frac{1}{N_{\text{traj}}(N_{\text{traj}}-1)} \sum_{i \ne j} \Theta(\epsilon - \|\mathbf{x}_i - \mathbf{x}_j\|_2)
      \end{equation*}
  \end{itemize}
  
  \textbf{Attention Jacobian Power Iteration:}
  \begin{itemize}
    \item Dense Attention Jacobian $\mathbf{J}_{\text{Attn}} \in \mathbb{R}^{(Nd) \times (Nd)}$ is memory-prohibitive.
    \item Spectral Norm $\|\mathbf{J}_{\text{Attn}}\|_2$ computed via JVP/VJP Power Iteration:
      \begin{align*}
        \mathbf{u}_k &= \mathbf{J}_{\text{Attn}} \mathbf{v}_k = \texttt{jvp}(\text{Attn}, \mathbf{X}, \mathbf{v}_k) \\
        \mathbf{w}_k &= \mathbf{J}_{\text{Attn}}^T \mathbf{u}_k = \texttt{vjp}(\text{Attn}, \mathbf{X}, \mathbf{u}_k)
      \end{align*}
    \item Rayleigh quotient update: $\sigma_{\max}(\mathbf{J}_{\text{Attn}}) \approx \sqrt{\frac{\langle \mathbf{w}_k, \mathbf{v}_k \rangle}{\langle \mathbf{v}_k, \mathbf{v}_k \rangle}}$
  \end{itemize}
\end{frame}

% Slide 10A: Perturbation Propagation and Operator Jacobians
\begin{frame}{Perturbation Propagation and Operator Jacobians}
  \textbf{Linearized Perturbation Propagation:}
  For a small input perturbation $\bm{\delta}_0$, the output perturbation after an operator $f(\mathbf{x})$ is governed by the local Jacobian matrix $\mathbf{J}_f(\mathbf{x}) = \frac{\partial f(\mathbf{x})}{\partial \mathbf{x}}$:
  \begin{equation*}
    \bm{\delta}_1 = f(\mathbf{x} + \bm{\delta}_0) - f(\mathbf{x}) \approx \mathbf{J}_f(\mathbf{x}) \bm{\delta}_0
  \end{equation*}
  
  \vspace{0.3cm}
  
  \textbf{Operator-Specific Jacobian Dynamics:}
  \begin{itemize}
    \item \textbf{Self-Attention ($\mathbf{J}_{\text{Attn}}$)}: Spatially couples tokens, routing and spreading the perturbation across the sequence dimension.
    \item \textbf{MLP Block ($\mathbf{J}_{\text{MLP}}$)}: Operates token-wise; high-dimensional expansion causes isotropic amplification (stretching, typically $\|\mathbf{J}_{\text{MLP}}\|_2 > 1$).
    \item \textbf{Normalization ($\mathbf{J}_{\text{Norm}}$)}: Rescales activations to a constant radial bound, acting as a contracting projection (folding component).
  \end{itemize}
\end{frame}

%\section{Results: Perturbation Propagation (Layers)}

% Slide 12: MLP Static Weights
\begin{frame}{MLP Static Weights}
  \begin{columns}
    \begin{column}{0.32\textwidth}
      \textbf{SwiGLU MLP:}
      $$ \text{MLP}(\mathbf{x}) = W_{\text{down}} \left( \text{SiLU}(W_{\text{gate}} \mathbf{x}) \odot (W_{\text{up}} \mathbf{x}) \right) $$

      \vspace{0.15cm}

      \textbf{Scaled Frobenius Norms:}
      $$ \tilde{F}_{\text{gate}} = \frac{\|W_{\text{gate}}\|_F^2}{d}, \quad \tilde{F}_{\text{up}} = \frac{\|W_{\text{up}}\|_F^2}{d} $$

      \vspace{0.15cm}

      \begin{itemize}
        \item Norms are $>1$ at all layers.
        \item Indicates baseline capacity to expand perturbations.
      \end{itemize}
    \end{column}
    \begin{column}{0.66\textwidth}
      \centering
      \includegraphics[width=0.99\textwidth]{plots_new/deepseek/jacobians_MLP/scaled_frobenius_jacobian_base_All_Prompts_Aggregated_aggregated_harmonic_xscale-linear_yscale-log.png}
      \\[0.1cm]
      {\tiny Fig 2: Scaled Frobenius norms of MLP static weights (DeepSeek).}
    \end{column}
  \end{columns}
\end{frame}

% Slide 13: MLP Activation Densities
\begin{frame}{MLP Activation Densities}
  \begin{columns}
    \begin{column}{0.32\textwidth}
      \textbf{Gate Activation Density:}
      $$ \overline{s_x^2} = \frac{\|\text{SiLU}(W_{\text{gate}} \mathbf{x})\|_2^2}{d_{\text{ff}}} $$

      \vspace{0.15cm}

      \textbf{Derivative Density:}
      $$ \overline{d_x^2} = \frac{\|(W_{\text{up}} \mathbf{x}) \odot \text{SiLU}'(W_{\text{gate}} \mathbf{x})\|_2^2}{d_{\text{ff}}} $$

      \vspace{0.15cm}

      \begin{itemize}
        \item Drives dynamic Jacobian stretching.
        \item Grows significantly in deeper layers.
      \end{itemize}
    \end{column}
    \begin{column}{0.66\textwidth}
      \centering
      \includegraphics[width=0.99\textwidth]{plots_new/deepseek/jacobians_MLP/activation_densities_jacobian_base_All_Prompts_Aggregated_aggregated_harmonic_xscale-linear_yscale-log.png}
      \\[0.1cm]
      {\tiny Fig 3: Dynamic activation densities across layers (DeepSeek).}
    \end{column}
  \end{columns}
\end{frame}

% Slide 14: MLP Jacobian Norm and Isotropic Expansion
\begin{frame}{MLP Jacobian Norm and Isotropic Expansion}
  \begin{columns}
    \begin{column}{0.32\textwidth}
      \textbf{Maximum Expansion (Spectral Norm):} 
      $$\|\mathbf{J}(\mathbf{x})\|_2 = \sigma_{\max}(\mathbf{J}(\mathbf{x}))$$

      \vspace{0.15cm}

      \textbf{Expected Isotropic Expansion:} 
      $$\bar{\lambda} = \frac{1}{d} \|\mathbf{J}(\mathbf{x})\|_F^2$$

      \vspace{0.15cm}

      \begin{itemize}
        \item $\|\mathbf{J}(\mathbf{x})\|_2 > 1$ at all layers.
        \item $\bar{\lambda} > 1$ (except first few layers).
      \end{itemize}
    \end{column}
    \begin{column}{0.66\textwidth}
      \centering
      \includegraphics[width=0.99\textwidth]{plots_new/deepseek/jacobians_MLP/jacobian_together_jacobian_base_All_Prompts_Aggregated_aggregated_harmonic_xscale-linear_yscale-log.png}
      \\[0.1cm]
      {\tiny Fig 4: Jacobian spectral norm and isotropic expansion (DeepSeek).}
    \end{column}
  \end{columns}
\end{frame}

% Slide 15: MLP Mean-Field Theory (MFT)
\begin{frame}{MLP Mean-Field Theory (MFT)}
  \textbf{MFT Simplification Strategy (Vanilla MLP $\mathbf{J} = W_2 \Sigma W_1$):}
  \begin{equation*}
    \|\mathbf{J}\|_F^2 = \text{Tr}(H_2 \Sigma H_1 \Sigma) = \sum_{i,j} (H_2)_{ij} (H_1)_{ji} \phi'_i \phi'_j
  \end{equation*}
  
  \textbf{Core Assumptions:}
  \begin{enumerate}
    \item Neglect off-diagonal terms: $\sum_{i \neq j} (H_2)_{ij} (H_1)_{ji} \phi'_i \phi'_j \approx 0$
    \item Row norms of weight matrices are uniform: $(H_1)_{ii} \approx \frac{\|W_1\|_F^2}{d_{\text{ff}}}$
  \end{enumerate}

  \textbf{Closed-Form Isotropic Expansion Estimate:}
  \begin{equation*}
    \overline{\lambda} \approx \left( \frac{\|W_1\|_F^2}{d} \right) \left( \frac{\|W_2\|_F^2}{d_{\text{ff}}} \right) \gamma, \quad \gamma = \frac{1}{d_{\text{ff}}} \sum_{i=1}^{d_{\text{ff}}} (\phi'_i)^2
  \end{equation*}
  \begin{itemize}
    \item Predicts the expansion norm without computing dense Jacobians.
  \end{itemize}
\end{frame}

% Slide 16A: MFT Assumption Validation
\begin{frame}{MFT Assumption Validation}
  \begin{columns}
    \begin{column}{0.30\textwidth}
      \textbf{Assumption Validation:}
      \begin{itemize}
        \item Diagonal dominance ($R \approx 1$) and Uniformity ($CV \ll 1$) are well-satisfied, showing slight discrepancies only in the last layers.
      \end{itemize}
    \end{column}
    \begin{column}{0.68\textwidth}
      \centering
      \includegraphics[width=0.48\textwidth]{plots_new/deepseek/mft-check/mft_validation_plot_assumption1_jacobian_base_All_Prompts_Aggregated_aggregated.png}
      \includegraphics[width=0.48\textwidth]{plots_new/deepseek/mft-check/mft_validation_plot_assumption2_jacobian_base_All_Prompts_Aggregated_aggregated.png}
      \\[0.1cm]
      {\tiny Fig 5a: MFT assumptions: diagonal dominance (left) and uniformity (right).}
    \end{column}
  \end{columns}
\end{frame}

% Slide 16B: MFT Prediction Verification
\begin{frame}{MFT Prediction Verification}
  \begin{columns}
    \begin{column}{0.30\textwidth}
      \textbf{Isotropic Expansion Prediction:}
      \begin{itemize}
        \item Outliers in layer with large Jacobian using mean 
        \item mean is sensitive to outliers 
        \item The 25-75 percentile bands agree consistently
      \end{itemize}
    \end{column}
    \begin{column}{0.68\textwidth}
      \centering
      \includegraphics[width=0.99\textwidth]{plots_new/deepseek/jacobians_MLP/mft_comparison_jacobian_base_All_Prompts_Aggregated_aggregated_xscale-linear_yscale-log.png}
      \\[0.1cm]
      {\tiny Fig 5b: Measured isotropic expansion vs. closed-form MFT prediction (DeepSeek).}
    \end{column}
  \end{columns}
\end{frame}





% Slide 17: Attention Static Weights
\begin{frame}{Attention Static Weights}
  \begin{columns}
    \begin{column}{0.32\textwidth}
      \textbf{Routing Matrix ($W_{QK}$):}
      $$W_{QK} = \frac{1}{\sqrt{d_k}} W_Q W_K^T$$

      \vspace{0.15cm}

      \textbf{Mixing Matrix ($W_{VO}$):}
      $$W_{VO} = W_V W_O$$

      \vspace{0.15cm}

      \begin{itemize}
        \item Routing and mixing norms are consistently $> 1$ at all layers.
        \item Confirms capacity to stretch perturbations.
      \end{itemize}
    \end{column}
    \begin{column}{0.66\textwidth}
      \centering
      \includegraphics[width=0.99\textwidth]{plots_new/deepseek/jacobians_attn/attn_static_weights_jacobian_base_aggregated_harmonic_xscale-linear_yscale-log.png}
      \\[0.1cm]
      {\tiny Fig 6: Spectral norms of routing and mixing weight matrices (DeepSeek).}
    \end{column}
  \end{columns}
\end{frame}

% Slide 18: Attention Mixing and Routing
\begin{frame}{Attention Mixing and Routing}
  \textbf{Self-Attention Output Formulation:}
  \begin{equation*}
    \mathbf{Y} = \mathbf{A} \mathbf{X} W_V W_O, \quad \delta \mathbf{Y} = \underbrace{(\delta \mathbf{A}) \mathbf{X} W_V W_O}_{\text{Nonlinear Routing}} + \underbrace{\mathbf{A} (\delta \mathbf{X}) W_V W_O}_{\text{Static Linear Mixing}}
  \end{equation*}

  \textbf{Static Linear Mixing Bounds:}
  \begin{equation*}
    \left\| (\mathbf{A} \delta \mathbf{V})_i \right\|_2 \le \max_j \|\delta \mathbf{v}_j\|_2 \le \left(\max_j \|\delta \mathbf{x}_j\|_2\right) \left\| W_V W_O \right\|_2
  \end{equation*}
  \begin{itemize}
    \item Softmax row-stochasticity ($\sum_j A_{i,j}=1$) ensures linear mixing itself cannot autonomously amplify the perturbation.
  \end{itemize}

  \textbf{Nonlinear Routing Bounds:}
  \begin{equation*}
    \left\| \delta \mathbf{Y}_{\text{route}} \right\|_2 \le \frac{\|\mathbf{X}\|_2^2}{\sqrt{d_k}} \|\delta \mathbf{X}\|_2 \left\|W_V W_O\right\|_2 \left\|W_Q W_K^T\right\|_2
  \end{equation*}
  \begin{itemize}
    \item Bilinear coupling: stretching scales quadratically with sequence norm $\|\mathbf{X}\|_2$.
    \item Max softmax sensitivity bound $\|\mathbf{J}_{\text{sm}}(\mathbf{A})\|_2 \le 0.5$ (Popoviciu's inequality).
  \end{itemize}
\end{frame}

% % Slide 19: Attention Jacobian Spectral Norm & Entropy
% \begin{frame}{Attention Jacobian and Entropy}
%   \begin{columns}
%     \begin{column}{0.32\textwidth}
%       \textbf{Jacobian Scaling:}
%       \begin{itemize}
%         \item Global attention Jacobian norm $\|\mathbf{J}_{\text{Attn}}\|_2$ scales with $\sqrt{N}$ due to random matrix coupling.
%       \end{itemize}
      
%       \textbf{Softmax Entropy Ratio:}
%       \begin{itemize}
%         \item Softmax saturation could collapse routing sensitivity ($\mathbf{J}_{\text{sm}} \to \mathbf{0}$).
%         \item Empirically, normalized entropy ratio remains high ($0.4$--$0.8$), preventing Jacobian vanishing.
%       \end{itemize}
%     \end{column}
%     \begin{column}{0.66\textwidth}
%       \centering
%       \includegraphics[width=0.49\textwidth]{plots_new/deepseek/jacobians_attn/attn_spectral_norms_jacobian_base_aggregated_harmonic_xscale-linear_yscale-log.png}
%       \includegraphics[width=0.49\textwidth]{plots_new/microsoft/jacobians_attn/attn_spectral_norms_jacobian_base_aggregated_harmonic_xscale-linear_yscale-log.png}
%       \\[0.1cm]
%       {\tiny Fig 7: Global attention Jacobian spectral norm $\|\mathbf{J}_{\text{Attn}}\|_2$ (DeepSeek left, Phi-4 right).}
%     \end{column}
%   \end{columns}
% \end{frame}

% Slide 20: Normalization and Residual Connections
\begin{frame}{Normalization and Residual Connections}
  \textbf{Normalization (RMSNorm / LayerNorm) Jacobian:}
  \begin{equation*}
    \mathbf{J}_{\text{Norm}} = \frac{1}{S} \mathbf{P}_{\perp}
  \end{equation*}
  \begin{itemize}
    \item \textbf{Orthogonal Projection} ($\mathbf{P}_{\perp}$): Eliminates perturbations pointing in the radial direction of the state vector.
    \item \textbf{Radial Scaling} ($S^{-1}$): Inversely scales incoming perturbations by state magnitude/variance, counteracting exponential growth.
  \end{itemize}

  \textbf{Residual Connections:}
  \begin{equation*}
    \mathbf{J}_{\text{res}} = \mathbf{I} + \mathbf{J}_F
  \end{equation*}
  \begin{itemize}
    \item Shifts eigenvalues of the sub-block Jacobian ($\lambda_i \to 1+\lambda_i$).
    \item Preserves the radial components annihilated by $\mathbf{P}_{\perp}$.
    \item Interplay: Normalization folding + Attention/MLP stretching keeps the residual stream bounded but chaotic.
  \end{itemize}
\end{frame}

% Slide 21: Finite Size Perturbation Propagation
\begin{frame}{Finite Size Perturbation Propagation}
  \begin{columns}
    \begin{column}{0.4\textwidth}
      \textbf{Layer-wise Evolution:}
      \begin{itemize}
        \item A finite-size perturbation is applied.
        \item Perturbation magnitude increases rapidly in early layers.
        \item Plateaus in middle layers.
        \item Decays in the final layers (converging to output probability distribution).
      \end{itemize}
    \end{column}
    \begin{column}{0.58\textwidth}
      \centering
      \includegraphics[width=0.99\textwidth]{plots_new/deepseek/perturbations/single_token_perturb_all_All_Prompts_Aggregated_aggregated_metric-harmonic_err-fan_xscale-linear_yscale-log.png}
      \\[0.1cm]
      {\tiny Fig 8: Finite perturbation magnitude across model layers (DeepSeek).}
    \end{column}
  \end{columns}
\end{frame}

% Slide 22: Scaling and FP16 Precision Limits
\begin{frame}{Scaling and FP16 Precision Limits}
  \begin{columns}
    \begin{column}{0.30\textwidth}
      \textbf{Linear Scaling Law:}
      \begin{equation*}
        \log \|\delta \mathbf{x}_\ell\|_2 \approx \log A(\ell) + 1 \cdot \log r
      \end{equation*}
      \begin{itemize}
        \item Exponents close to 1 ($0.993$ and $0.986$) confirm tangent map linearity for medium perturbations.
      \end{itemize}
      
      \textbf{Numerical Precision Floor:}
      \begin{itemize}
        \item At small perturbation radii, scaling breaks due to FP16 quantization limit:
          \begin{equation*}
            \Delta x = 2^{-11} \approx 4.88 \times 10^{-4}
          \end{equation*}
      \end{itemize}
    \end{column}
    \begin{column}{0.68\textwidth}
      \centering
      \includegraphics[width=0.8\textwidth]{plots_new/deepseek/perturbations/single_token_perturb_all_All_Prompts_Aggregated_aggregated_scaling_rainbow_median_xscale-log_yscale-log.png} \\
      \includegraphics[width=0.8\textwidth]{plots_new/deepseek/perturbations/single_token_perturb_all_All_Prompts_Aggregated_aggregated_extracted_jacobians_together_harmonic_xscale-linear_yscale-log.png}
      \\[0.1cm]
      {\tiny Fig 9: Exponent scaling and linearized perturbation gain (DeepSeek).}
    \end{column}
  \end{columns}
\end{frame}

%\section{Results: Perturbation Propagation (Tokens)}

% Slide 23: Trajectory Divergence
\begin{frame}{Trajectory Divergence}
  \begin{columns}
    \begin{column}{0.32\textwidth}
      \textbf{Latent Perturbations:}
      \begin{itemize}
        \item Small prompt perturbations don't trigger immediate changes.
        \item Perturbation remains latent in the hidden states (identical tokens are selected).
      \end{itemize}
      
      \textbf{Token Flip Divergence:}
      \begin{itemize}
        \item At a critical step, the accumulated perturbation shifts logits past the softmax/sampling threshold.
        \item A token flip occurs, causing the trajectories to diverge rapidly.
      \end{itemize}
    \end{column}
    \begin{column}{0.66\textwidth}
      \centering
      \includegraphics[width=0.49\textwidth]{plots_new/deepseek/divergence/cos_timeseries_0_1_window_size_16.png}
      \includegraphics[width=0.49\textwidth]{plots_new/deepseek/divergence/cos_timeseries_0_1_window_size_16_sentence.png}
      \\[0.1cm]
      {\tiny Fig 10: Trajectory cosine distance of hidden states (left) vs. text embeddings (right).}
    \end{column}
  \end{columns}
\end{frame}

% Slide 24: Divergence Metrics Comparison
\begin{frame}{Divergence Metrics Comparison}
  \begin{columns}
    \begin{column}{0.30\textwidth}
      \textbf{Metric Comparison:}
      \begin{itemize}
        \item Comparison of different trajectory-wise distance measures (sliding window=16).
        \item Cross-correlation, DTW, Fréchet, and Hausdorff all show similar sharp step-like divergence.
        \item Indicates the qualitative results are robust and not an artifact of a specific metric.
      \end{itemize}
    \end{column}
    \begin{column}{0.68\textwidth}
      \centering
      \includegraphics[width=0.48\textwidth]{figures_pdf/divergence/childhood_personality_development_0_0.0004/hidden_states_layer_-1_sentence/20260213-161650-81b984/plots/cross_corr_pearson_0_2_window_size_16.pdf}
      \includegraphics[width=0.48\textwidth]{figures_pdf/divergence/childhood_personality_development_0_0.0004/hidden_states_layer_-1_sentence/20260213-161650-81b984/plots/dtw_timeseries_0_2_window_size_16.pdf} \\
      \includegraphics[width=0.48\textwidth]{figures_pdf/divergence/childhood_personality_development_0_0.0004/hidden_states_layer_-1_sentence/20260213-161650-81b984/plots/frechet_timeseries_0_2_window_size_16.pdf}
      \includegraphics[width=0.48\textwidth]{figures_pdf/divergence/childhood_personality_development_0_0.0004/hidden_states_layer_-1_sentence/20260213-161650-81b984/plots/hausdorff_timeseries_0_2_window_size_16.pdf}
      \\[0.1cm]
      {\tiny Fig 11: Divergence comparison across four metrics.}
    \end{column}
  \end{columns}
\end{frame}


%\section{Results: Recurrence and Dimensionality}

% Slide 26: Recurrence Plots: Layer Comparison
\begin{frame}{Recurrence Plots: Layer Comparison}
  \begin{columns}
    \begin{column}{0.30\textwidth}
      \textbf{Layer Differences:}
      \begin{itemize}
        \item RPs generated at fixed recurrence rate $\text{RR}=0.03$.
        \item First Layer: Lacks coherent structures (shows embedding patterns).
        \item Last Layer: Clear diagonal line structures emerge.
      \end{itemize}
      \textbf{Significance:}
      \begin{itemize}
        \item Fragmented diagonals reveal nearby trajectories remaining close for short intervals before separating.
        \item Classic signature of chaotic attractor dynamics.
      \end{itemize}
    \end{column}
    \begin{column}{0.68\textwidth}
      \centering
      \includegraphics[width=0.48\textwidth]{figures_pdf/RP/compare_first_last/interstellar_propulsion_review_recurrence_first_rr_0.03.pdf}
      \includegraphics[width=0.48\textwidth]{figures_pdf/RP/compare_first_last/interstellar_propulsion_review_recurrence_last_rr_0.03.pdf} \\
      \includegraphics[width=0.48\textwidth]{figures_pdf/RP/compare_first_last/quantum_consciousness_hallucination_recurrence_first_rr_0.03.pdf}
      \includegraphics[width=0.48\textwidth]{figures_pdf/RP/compare_first_last/quantum_consciousness_hallucination_recurrence_last_rr_0.03.pdf}
      \\[0.1cm]
      {\tiny Fig 13: Recurrence plots of first vs. last layer.}
    \end{column}
  \end{columns}
\end{frame}

% Slide 27: Recurrence Plots: Semantic Features
\begin{frame}{Recurrence Plots: Semantic Features}
  \begin{columns}
    \begin{column}{0.30\textwidth}
      \textbf{Semantic Signatures in RPs:}
      \begin{itemize}
        \item \textbf{Chain-of-Thought Block}: Identified as a dark block at bottom left, indicating a stable "reasoning state".
        \item \textbf{Templates/Lists}: Quasi-periodic grid-like patterns indicate repetitive output structure.
        \item \textbf{Edge of Chaos}: Short interrupted diagonals correspond to chaotic transitions.
      \end{itemize}
    \end{column}
    \begin{column}{0.68\textwidth}
      \centering
      \includegraphics[width=0.48\textwidth]{figures_pdf/RP/examples_llm/recurrence_traj_2_rr_0.05.pdf}
      \includegraphics[width=0.48\textwidth]{figures_pdf/RP/examples_llm/recurrence_traj_3_rr_0.05.pdf} \\
      \includegraphics[width=0.48\textwidth]{figures_pdf/RP/examples_llm/recurrence_traj_0_rr_0.05.pdf}
      \includegraphics[width=0.48\textwidth]{figures_pdf/RP/examples_llm/recurrence_traj_6_rr_0.05.pdf}
      \\[0.1cm]
      {\tiny Fig 14: Last-layer RPs for different prompts (RR=0.05).}
    \end{column}
  \end{columns}
\end{frame}


% Slide 28A: Correlation Dimension
\begin{frame}{Correlation Dimension}
  % The Correlation Dimension $D_2$ is defined via the correlation sum $C(\epsilon)$:
  \begin{equation*}
    C(\epsilon) = \frac{1}{N(N-1)} \sum_{i \ne j} \Theta(\epsilon - \|\mathbf{x}_i - \mathbf{x}_j\|_2) % \quad \implies \quad C(\epsilon) \propto \epsilon^{D_2} \quad \text{as } \epsilon \to 0
  \end{equation*}

  Correlation dimension allows one to distinguish between deterministic chaos and random noise

  Not enough datapoints to compute exact dimension.

  
  
  \vspace{0.3cm}
  
  \centering
  \begin{tabular}{ccc}
    \includegraphics[width=0.31\textwidth]{plots/dimension/random.png} &
    \includegraphics[width=0.31\textwidth]{plots/dimension/childhood_personality_development/cosine_sim_first_first.png} &
    \includegraphics[width=0.31\textwidth]{plots/dimension/childhood_personality_development/cosine_sim_last_last.png} \\
    \small{RANDOM} & \small{FIRST LAYER} & \small{LAST LAYERS}
  \end{tabular}
\end{frame}

% Slide 28B: Distance Distribution
\begin{frame}{Distance Distribution}
  The internal states at the first layers resemble a set of random vectors.
  
  \vspace{0.5cm}
  
  \centering
  \begin{tabular}{ccc}
    \includegraphics[width=0.31\textwidth]{plots/dimension/random_dist.png} &
    \includegraphics[width=0.31\textwidth]{plots/dimension/childhood_personality_development/cosine_sim_first_first_dist.png} &
    \includegraphics[width=0.31\textwidth]{plots/dimension/childhood_personality_development/cosine_sim_last_last_dist.png} \\
    \small\textbf{RANDOM} & \small\textbf{FIRST LAYER} & \small\textbf{LAST LAYERS}
  \end{tabular}
\end{frame}

% Slide 32: Conclusions
\begin{frame}{Conclusions}
  \begin{itemize}
    \item LLMs exhibit chaotic behavior especially in the last layers: Sensitivity to initial conditions and bounded dynamics
    \item Correlation dimension indicates deterministic chaos.
    \item Recurrence Plots and RQA suggest similarity to other chaotic systems.
    \item Identified stretching components (MLP expansion, Attention) balanced by folding components (normalization and residual connections).
    \item Derived and validated a closed-form Mean-Field approximation of MLP stretching.
    
  \end{itemize}
\end{frame}

% Slide 36: Acknowledgements
\begin{frame}{Acknowledgements}
  \begin{itemize}
    \item \textbf{Supervisors}: Special thanks to Dr. János Török and Kristóf Benedek for guidance.
    \item \textbf{Funding}: Supported by the National Research, Development and Innovation Fund and the Ministry of Culture and Innovation of Hungary (TKP2021-NVA-02 or TKP2021-EGA-02).
    \item \textbf{Compute}: Calculations performed on the HPC cluster of BME Department of Theoretical Physics.
  \end{itemize}
  \vfill
  \centering
  \textbf{Thank you for your attention!} \\
  Questions?
\end{frame}

% Slide 34: Memory-Efficient Jacobian Spectral Norm
\begin{frame}{Memory-Efficient Jacobian Spectral Norm}
  \begin{columns}
    \begin{column}{0.48\textwidth}
        Spectral Norm (largest singular value):
          \begin{equation*}
            \|J\|_2 = \sigma_{\max}(J) = \sqrt{\lambda_{\max}(J^T J)}
          \end{equation*}
      
      \textbf{Memory Bottleneck:}

      Full Jacobian size ($\approx 618$ GB in FP32).

      Forward pass on GPU; clone and copy hidden states to CPU immediately.

      Compute Jacobians layer-by-layer by loading only one layer's tensors to GPU.
    \end{column}
    \begin{column}{0.48\textwidth}
      \textbf{Matrix-Free Power Iteration:}
      % \begin{itemize}
      %   \item Find $\lambda_{\max}(J^T J)$ without materializing $J$ or $J^T J$:
      % \end{itemize}
      random unit vector $v$.

      Iterate for $K$ steps:

      \textbf{JVP}:
        \begin{equation*}
          u = J v \quad \text{via} \quad \texttt{torch.func.jvp}
        \end{equation*}

      \textbf{VJP}:
        \begin{equation*}
          w = J^T u \quad \text{via} \quad \texttt{torch.func.vjp}
        \end{equation*}
        (using reverse-mode autodiff)

      \textbf{Re-normalize}:

      Output spectral norm:
          \begin{equation*}
            \sigma_{\max} = \sqrt{\frac{v^T w}{v^T v}}
          \end{equation*}
    \end{column}
  \end{columns}
\end{frame}

% Slide 35: The Correlation Dimension (Grassberger & Procaccia)
\begin{frame}{Correlation Dimension (Grassberger \& Procaccia)}
  \begin{equation*}
        C(\epsilon) = \lim_{N_{\text{traj}} \to \infty} \frac{1}{N_{\text{traj}}(N_{\text{traj}}-1)} \sum_{i \ne j} \Theta(\epsilon - \|\mathbf{x}_i - \mathbf{x}_j\|_2)
      \end{equation*}
  \begin{quote}
    \scriptsize
    ``We study the correlation exponent $\nu$ introduced recently as a characteristic measure of strange attractors which allows one to distinguish between deterministic chaos and random noise. The exponent $\nu$ is closely related to the fractal dimension and the information dimension, but its computation is considerably easier. Its usefulness in characterizing experimental data which stem from very high dimensional systems is stressed.''
  \end{quote}
  
  \vfill
  \begin{flushright}
    \scriptsize
    Peter Grassberger and Itamar Procaccia (1983) \\
    \textbf{Measuring the Strangeness of Strange Attractors} \\
    \textit{Physica D: Nonlinear Phenomena}
  \end{flushright}
\end{frame}

\end{document}
